\section{Simultaneous Gaussian warnings and integrated activities}
\label{sec:gaussian-warnings}

The expansion contains factors which are small only after integration.
A warning factor, for example, need not be small on a configuration
with a large fluctuation. What is small is the probability that all
warnings in a given term are present. Their fluctuations are correlated,
and several warnings may be witnessed by the same Gaussian coordinate.
We first account for both effects explicitly.

\subsection{A bound for correlated coordinates}
Let $I$ be a finite coordinate set, and let $Z$ be a centred real Gaussian
vector with covariance matrix $A$. Suppose
\begin{equation}\label{eq:warning-covariance}
 A\ge0,\qquad v^{\mathsf T}Av\le M\sum_{i\in I}v_i^2
 \quad(v\in\mathbb R^I),\qquad M\ge0.
\end{equation}
The covariance may be singular. For $J\subseteq I$, $t\in\mathbb R$
and $s\ge0$,
\begin{equation}\label{eq:warning-rectangle}
 \mathbb P\{|Z_i|\ge t\text{ for every }i\in J\}
 \le 2^{|J|}\exp\{(-st+Ms^2/2)|J|\}.
\end{equation}
Indeed, for each choice of signs $\epsilon_i\in\{-1,1\}$ on $J$, put
$v_i=\epsilon_i$ on $J$ and $v_i=0$ elsewhere. On the corresponding
orthant event, $v\cdot Z\ge t|J|$. The Gaussian exponential identity gives
\[
 \mathbb E e^{s v\cdot Z}
 =e^{s^2v^{\mathsf T}Av/2}
 \le e^{Ms^2|J|/2}.
\]
For completeness, write $Z=A^{1/2}G$ with $G$ a vector of independent
standard real Gaussians. In each coordinate, completing the square in
$(2\pi)^{-1/2}\int e^{hu-u^2/2}\,du$ gives $e^{h^2/2}$;
multiplication proves the identity, including zero eigenvalues of $A$.
Markov's inequality applied to $e^{s v\cdot Z}$ bounds one orthant by
$e^{(-st+Ms^2/2)|J|}$, and the union of the $2^{|J|}$ orthants contains
the event in \eqref{eq:warning-rectangle}. This also proves the formula
when $J$ is empty.

Let $\Lambda$ be the coarse periodic lattice, with its lattice distance.
Write
\[
 b_D=(2D+1)^4.
\]
Every ball of radius $D$ has at most $b_D$ sites, including on small
periods where some representatives coincide. Assign each Gaussian
coordinate a site $g(i)\in\Lambda$, with at most $p$ coordinates at any
site. Warning addresses form a finite set $K$, with site map
$o:K\to\Lambda$ and at most $m$ addresses at any site. For a set
$X\subseteq K$, define
\[
 \mathcal C(X;t,D)=
 \{z:\ \forall a\in X\ \exists i\in I,
       \ d(o(a),g(i))\le D,\quad |z_i|\ge t\}.
\]
These are measurable sets, being finite intersections of finite unions
of closed coordinate inequalities.

\begin{proposition}\label{prop:warning-crowding}
If $H\ge0$ and
\begin{equation}\label{eq:warning-threshold}
 M/2+H b_Dm\le t,
\end{equation}
then
\begin{equation}\label{eq:warning-joint}
 \mathbb P\{Z\in\mathcal C(X;t,D)\}
 \le (2b_Dp\,e^{-H})^{|X|}.
\end{equation}
\end{proposition}
\begin{proof}
For each warning in $X$, choose a witnessing coordinate within distance
$D$. There are at most $(b_Dp)^{|X|}$ such assignments. For one assignment,
let $J$ be the set of distinct coordinates chosen. Then $|J|\le|X|$.
Conversely, any fixed chosen coordinate can witness at most $b_Dm$
addresses, since all such addresses lie within distance $D$ of its site.
Thus $|X|\le b_Dm|J|$. Applying \eqref{eq:warning-rectangle} with $s=1$,
the event belonging to this assignment has probability at most
\[
 2^{|J|}e^{(-t+M/2)|J|}
 \le 2^{|X|}e^{-Hb_Dm|J|}
 \le 2^{|X|}e^{-H|X|}.
\]
The union bound over assignments proves the assertion. If a warning
has no possible coordinate, its event is empty; the empty set of
warnings has probability one. Thus no nonemptiness exception is needed.
\end{proof}

\subsection{The scales used for warnings}
Fix the block length $L$. Let $h=h(x)=\widehat\beta_f$ be the fine
reference coupling, so that $h=x+c_L$ for a constant $c_L$ depending
on the fixed group and block length. For sufficiently large $x$, put
\begin{equation}\label{eq:warning-scales}
 B=\lceil(\log x)^2\rceil,\qquad
 P=\exp\{(\log\log h)^2\},\qquad
 d=1+B+(\log h)^2,\qquad \tau=P/\sqrt h.
\end{equation}
The actual fine coupling satisfies $h/2\le\beta\le2h$.
The fluctuation is written
$\mu+\beta^{-1/2}\operatorname{erect}(Z)$; $\mu$ is its mean and
$\operatorname{erect}$ identifies colour coordinates with bond variables.
For a fixed $\xi>0$, the coordinate warning threshold is
\[
 t=\xi\tau\sqrt\beta\ge \xi P/2.
\]
The weaker constant $1/2$ is convenient and follows from
$\sqrt\beta\ge\sqrt h/2$.

For the actual warning alphabet, its density per coarse site is bounded
by the fixed integer
\[
 m=L^4\#\{\text{plaquette axes}\}
       +2\#\{\text{axes}\}+1+L^4
       +L^4\#\{\text{axes}\}.
\]
The five first types are fine plaquette masks, coarse bond masks,
coarse bond likelihoods, coarse site likelihoods, and fine site
contractions; the last type is a fine bond cutoff. The coordinate
density $p$ is the fixed bound for the tubular colour coordinates.
Both bounds count distinct addresses, including different types at the
same site.

Suppose that the physical covariance satisfies
\eqref{eq:warning-covariance} with a constant $M$ independent of the
period and of $x$. Set $D=2B$, $\varepsilon=\xi/4$, and
\begin{equation}\label{eq:warning-price}
 E(x)=2b_{2B}p\exp\{-\varepsilon\sqrt P\}.
\end{equation}
Then, for all sufficiently large $x$, uniformly in the period and
in $X$,
\begin{equation}\label{eq:warning-actual-price}
 \mathbb P\{Z\in\mathcal C(X;t,2B)\}\le E(x)^{|X|}.
\end{equation}
Here is the scale comparison needed to justify this assertion. First
$d\le7(\log h)^2$ eventually, since $h=x+c_L$, $\log x\le2\log h$
eventually, and $B<(\log x)^2+1$. For every fixed $a>0$, integer
$r\ge0$ and constant $C\ge0$,
\begin{equation}\label{eq:warning-monitor-pays}
 C d^r\le P^a\quad\hbox{eventually}.
\end{equation}
Taking logarithms, it suffices to compare a constant plus
$2r\log\log h$ with $a(\log\log h)^2$; the latter dominates.
Since $b_{2B}m$ is bounded by a fixed multiple of $d^4$,
\eqref{eq:warning-monitor-pays} implies
$b_{2B}m\le\sqrt P$ and $M/2\le\varepsilon P$ eventually. Consequently
\[
 M/2+\varepsilon\sqrt P\,b_{2B}m
 \le2\varepsilon P=\xi P/2\le t.
\]
Proposition~\ref{prop:warning-crowding}, with $H=\varepsilon\sqrt P$,
gives \eqref{eq:warning-actual-price}.

In fact $E$ decays faster than every negative power of $h$. To see this,
write $y=\log\log h$. Then $\log h=e^y$ and
$\sqrt P=e^{y^2/2}$. For each $N>0$,
\[
 \varepsilon e^{y^2/2}-N e^y-\log(2pb_{2B})\longrightarrow+\infty,
\]
because $\log b_{2B}=O(y)$ and $e^{y^2/2}/e^y\to\infty$.
Thus $E(x)\le h^{-N}$ eventually. The thresholds may depend on
$L,\xi,M,p,m,N$; they do not depend on the lattice period.

\subsection{Why a nonzero warning forces an excursion}
We spell out the deterministic step used to connect the preceding
probability estimate with the actual factors. Let $q$ be the number of
colour coordinates on one bond. On the event that all coordinates
within distance $B$ of a site obey $|Z_i|\le\xi\tau\sqrt\beta$,
\[
 \|\beta^{-1/2}\operatorname{erect}(Z)_b\|\le q\xi\tau,
 \qquad
 \|\mu_b+\beta^{-1/2}\operatorname{erect}(Z)_b\|
       \le (C+q\xi)\tau
\]
for the bonds in question, provided $\|\mu_b\|\le C\tau$ there.
The first inequality follows by the triangle inequality in an
orthonormal colour basis; it does not require a sharp Euclidean bound.
Choose $q\xi\le1$ and an outer cutoff plateau $a\ge C+1$.
The bond cutoff is then identically one throughout this smaller event.

The remaining masks use the corresponding fine and predictor
estimates. In the source construction these estimates supply, outside
the enlarged seed region, a local error bound $\delta\tau$ for each
fine or predictor reading, together with the guide bounds
$A\tau/L^2$ and $A\tau$, respectively. The local error has three
parts: the mean error, the change under the controlled fluctuation,
and the nonlinear remainder. Each is bounded by $\delta\tau/3$.
If the positive quadratic constants of the fine action and predictor
penalty are $c_f,c_p$, the resulting bounds are
\begin{equation}\label{eq:warning-plateau}
 \text{fine reading}\le \frac{2}{\sqrt{c_f}}(A/L^2+\delta)\tau,
 \qquad
 \text{predictor reading}\le \frac{2}{\sqrt{c_p}}(A+\delta)\tau.
\end{equation}
Choose $\delta=\sqrt{c_f}/32$, then $L$ large enough that
$A/L^2\le\delta$, and finally
$W>(2/\sqrt{c_p})(A+\delta)$. The first coefficient in
\eqref{eq:warning-plateau} is at most $1/8<1/4$, and the second is
strictly less than $W$. The corresponding fine and predictor masks
therefore equal one; the local likelihood factors also equal one by
the same local plateau estimates. Taking $\xi$ smaller if necessary
bounds the linear fluctuation error by $\delta\tau/3$; the two
remaining errors are eventually below $\delta\tau/3$ by their
strictly smaller scale estimates.

A warning factor is either a bounded cutoff times one of these
factors minus one, or a bond cutoff minus one. Factors inside their
compulsory seed region are defined to be zero. Hence, under the
stated local estimates, a nonzero warning at an address implies a
coordinate excursion within distance $2B$ of its assigned site.
The local estimates producing \eqref{eq:warning-plateau}, and the
uniform physical covariance bound, remain explicit analytic inputs
here. The probability and integration arguments do not assume or
replace their production.

\subsection{Integration before multiplication of prices}
Fix a tree sector. Its integration variables are a Haar variable $Y$
and the single Gaussian vector $Z$ above. Their reference measure
$\nu(dY)\gamma_A(dZ)$ is a product of probability measures. All factors
in a term are evaluated at the same $(Y,Z)$.

A term is indexed by a finite set $S$ of distinct vertices of four
types: connected packets, forced activities, exponential logarithmic
corrections, and warnings. Assume these factors are measurable and
satisfy the following uniform pointwise bounds:
\[
 \begin{array}{c|c}
 \text{vertex }i&|f_i(Y,Z)|\le c_i\\ \hline
 \text{packet with }n_i\text{ seeds}&
      C_0 e^{(s-g)n_i}\\
 \text{forced activity }i&v_i\\
 \text{logarithmic correction }i&w_i\\
 \text{warning }i&T
 \end{array}
\]
Here all displayed prices are nonnegative. The price $w_i$ bounds the
actual factor $e^{\ell_i}-1$ as well as the logarithmic quantity when
that additional estimate is used elsewhere. A disconnected packet is
assigned factor and price zero. Write $X(S)$ for the distinct warning
addresses in $S$. Suppose each individual nonzero warning forces its
coordinate excursion, as established under the local inputs above.
Thus for every finite $S$, a nonzero product forces
$Z\in\mathcal C(X(S);t,2B)$.
Then
\begin{equation}\label{eq:warning-sector-mass}
 \int\left|\prod_{i\in S}f_i(Y,Z)\right|\,d\nu\,d\gamma_A
 \le \prod_{i\in S}\widetilde c_i,
 \qquad
 \widetilde c_i=
 \begin{cases}E(x)T,&i\text{ a warning},\\c_i,&\text{otherwise}.
 \end{cases}
\end{equation}
Indeed the integrand is bounded by
$\mathbf1_{\mathcal C(X(S);t,2B)}(Z)\prod_{i\in S}c_i$.
It is integrable, being measurable and uniformly bounded on a
probability space. Fubini's theorem and
\eqref{eq:warning-actual-price} give a bound
$E(x)^{|X(S)|}\prod c_i$, which is exactly the right-hand side of
\eqref{eq:warning-sector-mass}. In particular, no product of marginal
warning probabilities has been used.

The same calculation covers one marked insertion. Let its distinguished
vertex be $i_*$ and let its measurable factor $F_*$ satisfy
$|F_*|\le w_*$. Multiply it by a cutoff taking values in $[0,1]$.
Define the marked term to be zero unless $i_*\in S$, and otherwise
replace $f_{i_*}$ by this cutoff times $F_*$. Applying
\eqref{eq:warning-sector-mass} to $S\setminus\{i_*\}$ gives
\begin{equation}\label{eq:warning-marked-mass}
 \int |\text{marked term}|\,d\nu\,d\gamma_A
 \le \mathbf1_{\{i_*\in S\}}w_*
       \prod_{i\in S\setminus\{i_*\}}\widetilde c_i.
\end{equation}
If the removed vertex was a warning, its probability price is removed
as well. The remaining warnings alone are used in the event bound.
This is the marked vertex price required for the connected sums of
Section~\ref{sec:mixed-activity-prices}.

For the actual factors one may take
$T=1+\#\{\text{frame tags}\}+\#\{\text{point tags}\}$.
Indeed the unmodified mask or likelihood lies in $[0,T]$, the enforcing
cutoff lies in $[0,1]$, and $T\ge1$, so its cutoff times the difference
from one has absolute value at most $T$. A bond cutoff minus one is
bounded by one. Thus the small warning price in
\eqref{eq:warning-sector-mass} is supplied by the proved Gaussian
estimate. The packet, forced, and logarithmic pointwise bounds in the
table are separate analytic obligations; \eqref{eq:warning-sector-mass}
is the precise passage from those bounds to integrated vertex products.
